\documentclass[11pt]{article}

\usepackage[letterpaper,margin=0.82in,headheight=15pt]{geometry}
\usepackage[T1]{fontenc}
\usepackage{lmodern}
\usepackage{microtype}
\usepackage{amsmath,amssymb,amsthm,bm,mathtools}
\usepackage{stmaryrd}
\usepackage{booktabs,longtable,tabularx,array,multirow}
\usepackage{enumitem}
\usepackage{xcolor}
\usepackage{fancyhdr}
\usepackage{graphicx}
\usepackage{tikz-cd}
\usepackage{hyperref}
\usepackage[nameinlink,noabbrev]{cleveref}

\definecolor{navy}{HTML}{17324D}
\definecolor{teal}{HTML}{147D92}
\definecolor{orange}{HTML}{D47A27}
\definecolor{softblue}{HTML}{EEF4F7}
\definecolor{softgray}{HTML}{F3F5F6}
\definecolor{darkgray}{HTML}{46545E}
\definecolor{softorange}{HTML}{FFF4E8}
\definecolor{softgreen}{HTML}{EDF7F0}
\definecolor{softred}{HTML}{FCEEEE}

\hypersetup{
  colorlinks=true,
  linkcolor=navy,
  citecolor=teal,
  urlcolor=teal,
  pdftitle={Volume XIV: Continuum-Calibrated NNpi Dynamics, Hamiltonian-Generated Exchange Currents, and Nuclear Separator Flow},
  pdfauthor={DeuteronWigner project}
}

\pagestyle{fancy}
\fancyhf{}
\fancyhead[L]{\small\color{darkgray}DeuteronWigner formalism - Volume XIV}
\fancyhead[R]{\small\color{darkgray}Continuum NNpi and exchange-current completion}
\fancyfoot[L]{\small\color{darkgray}31 July 2026}
\fancyfoot[R]{\small\color{darkgray}\thepage}
\renewcommand{\headrulewidth}{0.4pt}

\setlist[itemize]{leftmargin=1.5em,itemsep=2pt,topsep=3pt}
\setlist[enumerate]{leftmargin=1.7em,itemsep=2pt,topsep=3pt}
\setlength{\parindent}{0pt}
\setlength{\parskip}{5pt}
\setlength{\emergencystretch}{2em}
\renewcommand{\arraystretch}{1.16}
\allowdisplaybreaks

\newtheorem{definition}{Definition}[section]
\newtheorem{axiom}[definition]{Axiom}
\newtheorem{principle}[definition]{Design principle}
\newtheorem{requirement}[definition]{Requirement}
\newtheorem{proposition}[definition]{Proposition}
\newtheorem{theorem}[definition]{Theorem}
\newtheorem{lemma}[definition]{Lemma}
\newtheorem{corollary}[definition]{Corollary}
\newtheorem{remark}[definition]{Remark}
\newtheorem{decision}[definition]{Model decision}

\crefname{definition}{definition}{definitions}
\crefname{axiom}{axiom}{axioms}
\crefname{principle}{design principle}{design principles}
\crefname{requirement}{requirement}{requirements}
\crefname{proposition}{proposition}{propositions}
\crefname{theorem}{theorem}{theorems}
\crefname{lemma}{lemma}{lemmas}
\crefname{corollary}{corollary}{corollaries}
\crefname{remark}{remark}{remarks}
\crefname{decision}{model decision}{model decisions}

\newcommand{\kT}{\bm{k}_{T}}
\newcommand{\pT}{\bm{p}_{T}}
\newcommand{\qT}{\bm{q}_{T}}
\newcommand{\lT}{\bm{\ell}_{T}}
\newcommand{\rT}{\bm{r}_{T}}
\newcommand{\sT}{\bm{s}_{T}}
\newcommand{\DT}{\bm{\Delta}_{T}}
\newcommand{\bD}{\bm{b}_{\Delta}}
\newcommand{\bM}{\bm{b}_{\mathrm{TMD}}}
\newcommand{\RT}{\bm{R}_{T}}
\newcommand{\dd}{\mathrm{d}}
\newcommand{\Tr}{\operatorname{Tr}}
\newcommand{\tr}{\operatorname{tr}}
\newcommand{\diag}{\operatorname{diag}}
\newcommand{\rank}{\operatorname{rank}}
\newcommand{\Span}{\operatorname{span}}
\newcommand{\Ker}{\operatorname{Ker}}
\newcommand{\Ran}{\operatorname{Ran}}
\newcommand{\Cov}{\operatorname{Cov}}
\newcommand{\Var}{\operatorname{Var}}
\newcommand{\Id}{\operatorname{Id}}
\newcommand{\Hom}{\operatorname{Hom}}
\newcommand{\End}{\operatorname{End}}
\newcommand{\Res}{\operatorname{Res}}
\newcommand{\PV}{\operatorname{PV}}
\newcommand{\Disc}{\operatorname{Disc}}
\newcommand{\Hil}{\mathcal H}
\newcommand{\LF}{\ensuremath{\mathrm{LF}}}
\newcommand{\TMD}{\ensuremath{\mathrm{TMD}}}
\newcommand{\GTMD}{\ensuremath{\mathrm{GTMD}}}
\newcommand{\GPD}{\ensuremath{\mathrm{GPD}}}
\newcommand{\PDF}{\ensuremath{\mathrm{PDF}}}
\newcommand{\QCD}{\ensuremath{\mathrm{QCD}}}
\newcommand{\EFT}{\ensuremath{\mathrm{EFT}}}
\newcommand{\TTN}{\ensuremath{\mathrm{TTN}}}
\newcommand{\TN}{\ensuremath{\mathrm{TN}}}
\newcommand{\PCAC}{\ensuremath{\mathrm{PCAC}}}
\newcommand{\cO}{\mathcal O}
\newcommand{\cM}{\mathcal M}
\newcommand{\cR}{\mathcal R}
\newcommand{\cS}{\mathcal S}
\newcommand{\cT}{\mathcal T}
\newcommand{\cV}{\mathcal V}
\newcommand{\cW}{\mathcal W}
\newcommand{\cE}{\mathcal E}
\newcommand{\cG}{\mathcal G}
\newcommand{\cH}{\mathcal H}
\newcommand{\cK}{\mathcal K}
\newcommand{\cQ}{\mathcal Q}
\newcommand{\cP}{\mathcal P}
\newcommand{\cC}{\mathcal C}
\newcommand{\cZ}{\mathcal Z}
\newcommand{\cD}{\mathcal D}
\newcommand{\cN}{\mathcal N}
\newcommand{\cU}{\mathcal U}
\newcommand{\cF}{\mathcal F}
\newcommand{\cA}{\mathcal A}
\newcommand{\cB}{\mathcal B}
\newcommand{\cL}{\mathcal L}
\newcommand{\cJ}{\mathcal J}
\newcommand{\cI}{\mathcal I}
\newcommand{\eps}{\varepsilon}
\newcommand{\abs}[1]{\left|#1\right|}
\newcommand{\norm}[1]{\left\lVert#1\right\rVert}
\newcommand{\inner}[2]{\left\langle #1\middle|#2\right\rangle}
\newcommand{\avg}[1]{\left\langle #1\right\rangle}
\newcommand{\phys}{\mathrm{phys}}
\newcommand{\eff}{\mathrm{eff}}
\newcommand{\bare}{\mathrm{bare}}
\newcommand{\ind}{\mathrm{ind}}
\newcommand{\ct}{\mathrm{ct}}
\newcommand{\reg}{\mathrm{reg}}
\newcommand{\trunc}{\mathrm{trunc}}
\newcommand{\ren}{\mathrm{ren}}
\newcommand{\num}{\mathrm{num}}
\newcommand{\UV}{\mathrm{UV}}
\newcommand{\IR}{\mathrm{IR}}
\newcommand{\COM}{\mathrm{CM}}
\newcommand{\Lz}{L_z}
\newcommand{\Jz}{J^z}
\newcommand{\Amp}{\ensuremath{\mathbf{Amp}}}
\newcommand{\Dens}{\ensuremath{\mathbf{Dens}}}
\newcommand{\Match}{\ensuremath{\mathbf{Match}}}
\newcommand{\Red}{\ensuremath{\mathbf{Red}}}
\newcommand{\Proc}{\ensuremath{\mathbf{Proc}}}
\newcommand{\Infer}{\ensuremath{\mathbf{Infer}}}
\newcommand{\HNN}{\Hil_{NN}}
\newcommand{\HNNpi}{\Hil_{NN\pi}}
\newcommand{\HNtwo}{\Hil_{\mathrm{N2}}}
\newcommand{\ContinuumChannel}{\texttt{NNPi\allowbreak Continuum\allowbreak Channel}}
\newcommand{\SpectralDensity}{\texttt{Continuum\allowbreak Spectral\allowbreak Density}}
\newcommand{\FiniteVolumeMap}{\texttt{FiniteVolume\allowbreak Spectral\allowbreak Map}}
\newcommand{\TransitionKernel}{\texttt{NNPi\allowbreak Transition\allowbreak Kernel}}
\newcommand{\PoleResidue}{\texttt{Deuteron\allowbreak PoleAnd\allowbreak Residue}}
\newcommand{\GaugedTerm}{\texttt{Gauged\allowbreak Hamiltonian\allowbreak Term}}
\newcommand{\ExchangeCurrent}{\texttt{Exchange\allowbreak Current\allowbreak Operator}}
\newcommand{\CurrentCertificate}{\texttt{CurrentBasis\allowbreak Completeness\allowbreak Certificate}}
\newcommand{\BlockContinuity}{\texttt{Block\allowbreak Continuity\allowbreak Ledger}}
\newcommand{\SeparatorTrajectory}{\texttt{Separator\allowbreak Trajectory}}
\newcommand{\ExplicitInduced}{\texttt{ExplicitInduced\allowbreak Pion\allowbreak Comparison}}
\newcommand{\ContinuumCoherent}{\texttt{Continuum\allowbreak Coherent\allowbreak Kernel}}
\newcommand{\StateBundle}{\texttt{N2\allowbreak State\allowbreak Bundle}}
\newcommand{\ProvComplex}{\texttt{N2\allowbreak Provenance2\allowbreak Complex}}
\newcommand{\AssumptionBundle}{\texttt{AssumptionBundle}}
\newcommand{\PredictionPlan}{\texttt{N2\allowbreak Prediction\allowbreak Plan}}

\newcolumntype{Y}{>{\raggedright\arraybackslash}X}
\newcolumntype{L}[1]{>{\raggedright\arraybackslash}p{#1}}

\title{\vspace{-1.0cm}
{\color{navy}\bfseries Volume XIV: Continuum-Calibrated \(NN\pi\) Dynamics,\\[2mm]
Hamiltonian-Generated Exchange Currents, and Nuclear Separator Flow}\\[4mm]
\large Spectral support, pole and residue control, blockwise current continuity,\\
continuum-aware pion operators, and stable explicit--induced matching}
\author{DeuteronWigner project}
\date{31 July 2026}

\begin{document}
\maketitle

\begin{center}
\begin{minipage}{0.94\textwidth}
\colorbox{softblue}{\parbox{0.96\textwidth}{
\textbf{Status and purpose.}
C16/N1 established a normalized, spin-resolved
\(NN\oplus NN\pi\) validation state with complete charge channels, diagonal
and transition recoil, nucleon-active and pion-active operators, executable
pion overlap subtraction, a finite current ledger, a helicity-resolved
coherent pilot, and order-resolved microscopic partonic parents.  Its
\(NN\pi\) sector, however, remained a finite analytic/discrete realization;
its current family was deliberately incomplete at the continuum nuclear
order; and its separator, pole, residue, and coherent inputs were not yet
calibrated as one flow.  The present volume is the normative formalism for
C17/N2.  It promotes the transition sector to a typed direct-integral
continuum with a controlled finite-volume approximation, derives the
self-energy, pole, residue, and cut ledgers, generates every retained current
from the same gauged Hamiltonian, proves continuity block by block, and
requires the internal-pion, exchange-pion, overlap, induced-Hamiltonian,
current, and partonic-operator descriptions to flow together under a declared
separator.  It remains a validation theory: no physical pion TMD, complete
continuum chiral EFT, physical nuclear shadowing, \(\Delta\Delta\), compact
six-quark, QCD matching, evolution, process, or inference readiness is
claimed.
}}
\end{minipage}
\end{center}

\begin{abstract}
The finite \(NN\oplus NN\pi\) light-front construction supplies the correct
sector algebra but does not by itself define continuum transition strength,
a pole residue, or a complete current basis.  We formulate an N2 theory in
which the \(NN\pi\) sector is a direct integral of charge-, helicity-,
partial-wave-, and regulator-resolved continuum channels.  The same
transition kernel defines a dispersive self-energy, its discontinuity, the
deuteron pole condition, the wave-function residue, and a convergent
finite-volume quadrature sequence.  Currents are obtained by gauging every
retained Hamiltonian term, including nonlocal regulators, momentum-dependent
kernels, sector-changing vertices, induced Feshbach interactions, and
counterterms.  A machine-readable completeness certificate maps each
Hamiltonian term to its charge, current, contact, regulator, and continuity
partners.  The continuity equation is tested separately in the \(NN\to NN\),
\(NN\leftrightarrow NN\pi\), and \(NN\pi\to NN\pi\) blocks and separately in
each charge channel.  A separator trajectory moves pion-like strength among
internal-nucleon, explicit exchange, overlap, and induced descriptions while
leaving matched observables invariant within a declared truncation law.  The
pion-active and transition partonic operators, energy--momentum operators,
and norm kernel are transformed with the Hamiltonian.  A continuum-aware
helicity-resolved coherent pilot and a derived completely positive reduction
are included without identifying either with physical shadowing or partonic
Wilson rescattering.  We give formal requirements, analytic benchmarks,
negative tests, tensor-network and continuum convergence obligations,
software interfaces, and the acceptance boundary to explicit
\(\Delta\Delta\), compact six-quark, and hidden-color sectors.
\end{abstract}

\tableofcontents

\section{Role, scope, and inherited N1 contracts}
\label{sec:scope}

\subsection{What C16/N1 established}

The inherited state is
\begin{equation}
 \Hil_{\mathrm{N1}}=\HNN\oplus\HNNpi,
 \qquad
 |D,\Lambda\rangle_{\mathrm{N1}}
 =|\Psi_{NN}^{\Lambda}\rangle+|\Psi_{NN\pi}^{\Lambda}\rangle,
 \label{eq:n1-state}
\end{equation}
with
\begin{equation}
 Z_{NN}+Z_{NN\pi}=1.
 \label{eq:n1-normalization}
\end{equation}
The state retains the complete
\(pn\pi^0\), \(pp\pi^-\), and \(nn\pi^+\) charge basis, exact spin-1
quantum numbers, three-body and transition recoil maps, diagonal and
transition partonic operators, a matched pion subtraction, and a
Hamiltonian-consistent finite current family.  These identities are immutable
inputs to N2.

\subsection{Why N1 is not yet the continuum theory}

A finite basis may verify Hermiticity, selection rules, current attachments,
and deterministic subtraction, but it does not automatically define the
continuum spectral measure or the on-shell discontinuity.  In particular,
replacing
\begin{equation}
 \frac{1}{E-H_Q+i0}
 \quad\text{by}\quad
 \frac{1}{E-H_Q+i\eps}
 \label{eq:epsilon-not-support}
\end{equation}
with a finite numerical \(\eps\) does not create physical support.  The
continuum threshold, measure, channel normalization, and map from a finite
sequence must be declared independently.

\begin{decision}[N2 transition]
N2 retains the N1 sector and operator algebra but replaces the finite
transition oracle, within the N2 validation root, by a continuum-calibrated
spectral theory and a controlled finite-volume/discretized-continuum route.
The N1 result remains an immutable benchmark and is never added to N2 as a
second physical mechanism.
\end{decision}

\subsection{Inherited typed distinctions}

The following remain noninterchangeable:
\begin{equation}
 \texttt{PARTONIC\_WILSON\_STAPLE}
 \neq
 \texttt{NUCLEAR\_COHERENT\_PROPAGATION}
 \neq
 \texttt{TAGGED\_FINAL\_STATE\_INTERACTION},
 \label{eq:rescattering-types}
\end{equation}
and
\begin{equation}
 \texttt{INTERNAL\_NUCLEON\_PION}
 \neq
 \texttt{EXPLICIT\_NUCLEAR\_NNPI}
 \neq
 \texttt{INDUCED\_PION\_OPERATOR}.
 \label{eq:pion-types}
\end{equation}
Array shape, identical dimensions, or similar low-energy behavior cannot
supply an implicit adapter.

\subsection{Explicit nonclaims}

N2 does not claim:
\begin{itemize}
\item a physical pion TMD or GTMD;
\item a complete continuum chiral nuclear EFT;
\item a full all-orders nuclear current basis;
\item a physical diffractive shadowing or Glauber calculation;
\item explicit \(\Delta\Delta\), six-quark, or hidden-color dynamics;
\item LF-to-QCD matching, Collins--Soper evolution, process factorization, or
      inference readiness.
\end{itemize}

\section{Continuum channel space and spectral normalization}
\label{sec:continuum}

\subsection{Direct-integral Hilbert space}

Let \(\alpha\) denote the complete discrete channel identity,
\begin{equation}
 \alpha=
 \left(
 c_\pi,I_{NN},S_{NN},\ell_\rho,\ell_\lambda,
 J^P,J^z,\Lambda,\text{regulator},\text{member}
 \right).
 \label{eq:channel-id}
\end{equation}
The continuum sector is the direct integral
\begin{equation}
 \HNNpi^{\mathrm{cont}}
 =
 \bigoplus_{\alpha}
 \int_{E_{\mathrm{th},\alpha}}^{\infty}\!{}^{\oplus}
 \dd\mu_\alpha(E)\,
 \Hil_{E\alpha},
 \qquad
 \dd\mu_\alpha(E)=\rho_\alpha(E)\dd E.
 \label{eq:direct-integral}
\end{equation}
Generalized eigenstates satisfy
\begin{equation}
 \langle E',\alpha'|E,\alpha\rangle
 =\delta_{\alpha'\alpha}\,\delta(E'-E)
 \label{eq:continuum-normalization}
\end{equation}
under the selected energy normalization.  Momentum, invariant-mass, or
phase-space normalizations are permitted only through a unitary typed adapter
that transforms the transition kernel and all measures together.

\begin{requirement}[Channel completeness]
Every continuum object must retain the charge channel, total charge, isospin,
parity, \(J^z\), pair spin/isospin, pion charge, Jacobi partial waves,
threshold, active/spectator assignment, regulator, normalization convention,
member, and assumption-plan identity.
\end{requirement}

\subsection{Threshold behavior}

For a three-body channel with effective relative orbital order \(\ell_\alpha\),
the spectral density near threshold is parameterized only after deriving the
phase-space behavior,
\begin{equation}
 \rho_\alpha(E)\,
 \abs{v_\alpha(E)}^2
 \sim
 (E-E_{\mathrm{th},\alpha})^{\nu_\alpha}
 \left[1+\cO(E-E_{\mathrm{th},\alpha})\right],
 \label{eq:threshold-law}
\end{equation}
where \(\nu_\alpha\) follows from the declared kinematics and partial waves.
A fitted noninteger exponent may be used only as a validation oracle with an
explicit discrepancy term; it cannot replace the threshold derivation.

\subsection{Spectral projectors}

The continuum resolution of the identity is
\begin{equation}
 Q_{\mathrm{cont}}
 =
 \sum_\alpha
 \int_{E_{\mathrm{th},\alpha}}^\infty
 \dd\mu_\alpha(E)
 |E,\alpha\rangle\langle E,\alpha|,
 \label{eq:spectral-projector}
\end{equation}
possibly supplemented by discrete \(NN\pi\)-like bound or pseudobound states,
which must be listed separately.  The spectral projector is part of the
Hamiltonian identity.  A change in \(\rho_\alpha\), threshold, regulator, or
normalization invalidates all cached self-energies, currents, and partonic
operators using it.

\section{Finite-volume and discretized-continuum maps}
\label{sec:finite-volume}

\subsection{Neutral finite spectral approximation}

A finite sequence is represented by
\begin{equation}
 \dd\mu_\alpha^{(L)}(E)
 =
 \sum_{n=1}^{N_L}
 w_{n\alpha}^{(L)}
 \delta\!\left(E-E_{n\alpha}^{(L)}\right)\dd E,
 \label{eq:finite-measure}
\end{equation}
with positive or signed weights according to the declared quadrature.  Its
identity contains the volume or discretization parameter \(L\), level count,
boundary conditions, threshold, quadrature rule, smearing parameter, and map
to the continuum measure.

\begin{requirement}[Terminology]
The object is called \FiniteVolumeMap.  It may be called a Luescher or
Lellouch--Luescher map only if the corresponding two-body quantization,
phase-shift, normalization, and matrix-element relations are implemented with
their hypotheses.  A generic spectral quadrature must not inherit that name.
\end{requirement}

\subsection{Weak convergence}

For a smooth compactly supported test function \(f\), require
\begin{equation}
 \sum_n w_{n\alpha}^{(L)}f(E_{n\alpha}^{(L)})
 \longrightarrow
 \int \dd\mu_\alpha(E)f(E).
 \label{eq:weak-convergence}
\end{equation}
Principal-value kernels require a subtraction-aware convergence criterion,
not merely weak convergence.  For example,
\begin{align}
 \PV\int\dd E\,
 \frac{F(E)}{E_0-E}
 ={}&
 \int\dd E\,
 \frac{F(E)-F(E_0)\chi(E)}{E_0-E}
 \notag\\
 &+F(E_0)
 \PV\int\dd E\,
 \frac{\chi(E)}{E_0-E},
 \label{eq:pv-subtraction}
\end{align}
where \(\chi\) is a declared local subtraction function.  The finite sequence
must converge for both terms separately.

\subsection{Continuum-map error budget}

The finite-to-continuum residual is decomposed as
\begin{equation}
 \delta_{\mathrm{FV}}
 =
 \delta_{N_L}
 +\delta_L
 +\delta_{\mathrm{quad}}
 +\delta_{\mathrm{smear}}
 +\delta_{\mathrm{th}}
 +\delta_{\mathrm{map}},
 \label{eq:fv-error}
\end{equation}
where none of the terms is absorbed into a physical transition width.

\section{Transition kernel, self-energy, pole, and residue}
\label{sec:self-energy}

\subsection{Hermitian transition family}

Let \(P\) project onto \(NN\) and \(Q\) onto the continuum \(NN\pi\) sector.
The transition kernel is
\begin{equation}
 V_{QP}(E):\HNN\longrightarrow\Hil_{E\alpha},
 \qquad
 V_{PQ}(E)=V_{QP}^{\dagger}(E).
 \label{eq:transition-kernel}
\end{equation}
Its matrix element carries charge, helicity, partial-wave, recoil, regulator,
separator, and member identity.  Emission and absorption are never fitted
independently.

\subsection{Feshbach self-energy}

The projected resolvent is
\begin{equation}
 G_P(z)
 =
 \frac{1}{z-H_{PP}-\Sigma(z)},
 \qquad
 \Sigma(z)=H_{PQ}\frac{1}{z-H_{QQ}}H_{QP}.
 \label{eq:projected-resolvent}
\end{equation}
In the continuum basis,
\begin{equation}
 \Sigma_{ab}(z)
 =
 \sum_\alpha
 \int_{E_{\mathrm{th},\alpha}}^\infty
 \dd\mu_\alpha(E)
 \frac{V_{a\alpha}(E)V_{\alpha b}(E)}{z-E}.
 \label{eq:self-energy-integral}
\end{equation}
For \(z=E_0+i0\),
\begin{align}
 \Sigma(E_0+i0)
 ={}&
 \PV\sum_\alpha\int\dd\mu_\alpha(E)
 \frac{V_{a\alpha}(E)V_{\alpha b}(E)}{E_0-E}
 \notag\\
 &-i\pi
 \sum_\alpha
 \rho_\alpha(E_0)V_{a\alpha}(E_0)V_{\alpha b}(E_0)
 \theta(E_0-E_{\mathrm{th},\alpha}).
 \label{eq:self-energy-disc}
\end{align}
Thus, absorption vanishes exactly below every open threshold.

\subsection{Pole condition}

For an isolated bound-state pole \(E_D\) below threshold, the effective
pole equation is
\begin{equation}
 \det\!
 \left[
 E_D-H_{PP}-\Sigma(E_D)
 \right]=0.
 \label{eq:pole-condition}
\end{equation}
For a one-dimensional reference subspace,
\begin{equation}
 E_D-E_{0}-\Sigma(E_D)=0.
 \label{eq:scalar-pole}
\end{equation}
The residue is
\begin{equation}
 Z_D
 =
 \left[
 1-\left.\frac{\partial\Sigma(E)}{\partial E}\right|_{E_D}
 \right]^{-1},
 \label{eq:pole-residue}
\end{equation}
under this resolvent convention.  In a multidimensional subspace, the residue
is a matrix on the pole eigenspace and must be projected with the corresponding
left/right eigenvectors.

\subsection{State normalization and norm kernel}

The continuum component is generated by
\begin{equation}
 \omega(E_D)
 =
 \frac{1}{E_D-H_{QQ}}H_{QP},
 \label{eq:wave-operator}
\end{equation}
and the effective-space norm kernel is
\begin{equation}
 N(E_D)=P\left[1+\omega^{\dagger}(E_D)\omega(E_D)\right]P.
 \label{eq:norm-kernel}
\end{equation}
The physical projected state is normalized with \(N\), not with the Euclidean
norm of the \(P\)-space coefficient vector.  The residue, norm-kernel, and
sector probability ledgers must agree within the declared truncation.

\begin{proposition}[No-width theorem for a closed finite spectrum]
If the declared spectrum is finite, all intermediate energies are off shell,
and no continuum or finite-volume spectral rule is supplied, then the
absorptive part generated by \cref{eq:projected-resolvent} is exactly zero.
A finite numerical \(\eps\) does not change this statement.
\end{proposition}

\section{Calibration, holdouts, and identifiability}
\label{sec:calibration}

\subsection{Restricted calibration set}

The first N2 calibration may use only:
\begin{enumerate}
\item the deuteron pole or mass condition;
\item one transition normalization or low-energy spectral condition;
\item one charge/isospin condition;
\item one Hamiltonian-owned current normalization.
\end{enumerate}
No named pion, tensor, or deuteron TMD receives an independent parameter.

\subsection{Frozen holdouts}

Before optimization, freeze at least:
\begin{itemize}
\item a second transition-energy point;
\item one spectral moment or pole residue;
\item one pion-active partonic moment;
\item one tensor transition observable;
\item one independent current component;
\item one angular-condition observable;
\item one coherent-amplitude quantity.
\end{itemize}
A failed holdout must revise a Hamiltonian term, current attachment, continuum
map, separator discrepancy, or operator basis.  It cannot be repaired by a
TMD-specific coefficient.

\subsection{Jacobian and null directions}

For observables \(\cO_i\) and parameters \(\theta_a\), define
\begin{equation}
 J_{ia}=\frac{\partial\cO_i}{\partial\theta_a}.
 \label{eq:calibration-jacobian}
\end{equation}
Singular values of \(J\), the Fisher matrix, and profile likelihoods are
recorded.  A null direction remains visible unless a new independent physical
condition is introduced; it may not be fixed by an unrelated TMD
normalization.

\section{Gauging the Hamiltonian}
\label{sec:gauging}

\subsection{External-field construction}

Introduce a gauge-covariant Hamiltonian family
\begin{equation}
 H[A]=H[0]+\int\dd^4x\,J^\mu(x)A_\mu(x)+\cO(A^2),
 \label{eq:gauged-hamiltonian}
\end{equation}
so that
\begin{equation}
 J^\mu(x)
 =
 -\left.\frac{\delta H[A]}{\delta A_\mu(x)}\right|_{A=0}
 \label{eq:current-functional}
\end{equation}
under the chosen sign convention.  Gauging the dynamical equation attaches the
external field to every propagator, vertex, exchange kernel, contact term,
regulator, and normalization factor.  This is the operational content of the
``gauging of equations'' construction \cite{KvinikhidzeBlankleider1999I,KvinikhidzeBlankleider1999II}.

\subsection{Nonlocal kernels and regulators}

For a momentum-dependent kernel \(V(p',p)\), minimal replacement in only the
external one-body momenta is insufficient.  The gauged kernel
\(V^\mu(p',p;q)\) must satisfy the finite-difference Ward identity
\begin{equation}
 q_\mu V^\mu(p',p;q)
 =
 Q_f V(p'-q,p)-V(p',p+q)Q_i
 \label{eq:kernel-ward}
\end{equation}
with the appropriate charge operators.  A nonlocal regulator
\(R_\Lambda(p',p)\) contributes its own gauged term.  Using different
regulators in the force and current generally destroys continuity.

\subsection{Hamiltonian-term ownership}

Write
\begin{equation}
 H_{N2}
 =
 \sum_t H_t,
 \qquad
 J^\mu_{N2}
 =
 \sum_t J_t^\mu+
 J_{\ct}^\mu+J_{\ind}^\mu.
 \label{eq:term-current-sum}
\end{equation}
Each term \(H_t\) has one of three certified statuses:
\begin{enumerate}
\item charged and accompanied by all required attachments;
\item neutral, with a proof-backed zero current;
\item unresolved, which blocks a completeness claim.
\end{enumerate}

\section{Declared-order exchange-current basis}
\label{sec:current-basis}

\subsection{Required operator families}

At the declared N2 order the current basis contains, when generated by the
retained Hamiltonian,
\begin{longtable}{L{0.30\textwidth}L{0.62\textwidth}}
\toprule
Current class & Physical and algebraic role\\
\midrule
\endhead
Nucleon one-body charge/current & Gauging of the nucleon kinetic and intrinsic current terms.\\
Pion-in-flight current & Photon insertion on the propagating exchanged or explicit pion.\\
\(NN\leftrightarrow NN\pi\) transition current & Gauging of sector-changing emission and absorption vertices.\\
Contact/seagull current & Simultaneous coupling required by derivative or gauged contact interactions.\\
Pair current & Negative-energy or pair-reduction partner of the retained interaction representation.\\
Recoil/retardation current & Energy- and momentum-dependent corrections associated with nonstatic kernels.\\
Momentum-dependent interaction current & Gauging of explicit momentum dependence in the potential.\\
Regulator-gauging current & Current generated by nonlocal cutoff functions and endpoint prescriptions.\\
Induced Feshbach current & Operator generated when the explicit continuum sector is eliminated.\\
Current counterterm & Resolution-dependent renormalization term fixed by declared current conditions.\\
Charge-density correction & Interaction-dependent correction to the density entering the continuity identity.\\
EMT interaction terms & Pair, seagull, pion-exchange, recoil, and regulator partners for \(T^{\mu\nu}\).\\
Axial/pseudoscalar companions & Matched axial, pseudoscalar, pion-pole, and contact terms where supported.\\
\bottomrule
\end{longtable}
Consistent one- and two-nucleon currents in chiral EFT display precisely this
logic: potentials, charge operators, pion exchange, short-range terms, and
regularization must be treated as one family
\cite{Pastore2009,Koelling2009,Koelling2011,Krebs2020,Pastore2011}.

\subsection{Completeness certificate}

The \CurrentCertificate\ is a bipartite incidence record between Hamiltonian
terms and current operators.  For each \(H_t\) it stores:
\begin{itemize}
\item charge assignment;
\item source and target sectors;
\item longitudinal and transverse attachments;
\item contact and regulator partners;
\item induced-operator partners;
\item counterterms;
\item continuity identities;
\item unresolved transverse low-energy constants.
\end{itemize}
The status
\texttt{N2\_DECLARED\_ORDER\_EXCHANGE\_CURRENT\_BASIS\_COMPLETE}
may be issued only when no unexplained gap remains.

\subsection{Energy-dependent currents}

If an effective interaction depends on the energy transfer, the ordinary
static continuity equation is modified by terms involving derivatives of the
kernel or charge operator.  The current identity and the pole-residue
normalization must therefore be derived in the same energy convention.  A
static current inserted into an energy-dependent Hamiltonian without the
corresponding derivative terms is inconsistent.

\section{Blockwise continuity and spin-1 current closure}
\label{sec:continuity}

\subsection{Sector-block continuity equation}

Write
\begin{equation}
 H=
 \begin{pmatrix}
 H_P & V_{PQ}\\
 V_{QP} & H_Q
 \end{pmatrix},
 \quad
 \rho=
 \begin{pmatrix}
 \rho_{PP} & \rho_{PQ}\\
 \rho_{QP} & \rho_{QQ}
 \end{pmatrix},
 \quad
 J^\mu=
 \begin{pmatrix}
 J_{PP}^\mu & J_{PQ}^\mu\\
 J_{QP}^\mu & J_{QQ}^\mu
 \end{pmatrix}.
 \label{eq:block-operators}
\end{equation}
The finite-resolution identity is
\begin{equation}
 q_\mu J^\mu=[H,\rho]+\delta_{\trunc}.
 \label{eq:continuity-master}
\end{equation}
It implies, for example,
\begin{align}
 q_\mu J_{PP}^\mu
 ={}&
 H_P\rho_{PP}+V_{PQ}\rho_{QP}
 -\rho_{PP}H_P-\rho_{PQ}V_{QP}
 +\delta_{PP},
 \label{eq:continuity-pp}\\
 q_\mu J_{QP}^\mu
 ={}&
 V_{QP}\rho_{PP}+H_Q\rho_{QP}
 -\rho_{QP}H_P-\rho_{QQ}V_{QP}
 +\delta_{QP}.
 \label{eq:continuity-qp}
\end{align}
All four sector blocks must close separately.

\subsection{Charge-channel resolution}

The continuity ledger is evaluated separately in
\begin{equation}
 pn\pi^0,
 \qquad
 pp\pi^-,
 \qquad
 nn\pi^+,
 \label{eq:charge-channels}
\end{equation}
and separately in each supported helicity and tensor block.  A cancellation
that occurs only after mixing charge channels or incompatible tensor
representations is not sufficient.

\subsection{Signed residual ledger}

The residual is decomposed as
\begin{align}
 \delta_{\mathrm{cont}}={}&
 \delta_{1N}+\delta_{\pi\mathrm{-flight}}
 +\delta_{\mathrm{trans}}+\delta_{\mathrm{seagull}}
 +\delta_{\mathrm{pair}}
 \notag\\
 &+\delta_{\mathrm{recoil}}
 +\delta_{\mathrm{mom}}
 +\delta_{\mathrm{reg}}
 +\delta_{\mathrm{Feshbach}}
 +\delta_{\ct}
 +\delta_{\mathrm{basis}}
 +\delta_{\mathrm{FV}}.
 \label{eq:continuity-ledger}
\end{align}
Removing any required nonzero contribution must produce a signed defect whose
origin is retained.

\subsection{Spin-1 current and angular condition}

The current must reproduce charge, magnetic, and quadrupole structures from a
common set of helicity amplitudes.  Independent extractions from different
current components and the light-front angular condition remain holdouts.  A
component-specific normalization is forbidden.

\subsection{Energy--momentum tensor}

The EMT family includes one-body, pair, seagull, pion-exchange, recoil, and
interaction terms.  Exchange corrections to deuteron gravitational form
factors provide a useful benchmark for the need to include these pieces
consistently \cite{HeZahed2024}.  The N2 implementation is still a
finite-resolution validation of that operator logic, not a physical EMT
prediction.

\section{Continuum-calibrated pion-active and transition operators}
\label{sec:pion-operators}

\subsection{Nuclear splitting amplitude}

The continuum-calibrated nuclear splitting amplitude is denoted
\begin{equation}
 \cS_{NN\pi/D}^{\Lambda'\Lambda;\alpha}
 (E,\{y_i,\bm p_{iT}\},\DT),
 \label{eq:continuum-splitting}
\end{equation}
where \(\alpha\) carries the complete continuum channel identity.  Its
normalization is fixed by the state and pole-residue conventions; it has no
independent deuteron-fit normalization.

\subsection{Pion-active parent}

The diagonal pion-active contribution is
\begin{align}
 W_{a/D}^{\pi\mathrm{-active}}
 ={}&
 \sum_{\alpha}
 \int\dd\mu_\alpha(E)
 \int[\dd\Gamma_3]
 \cS_{NN\pi/D}^{\alpha}
 \notag\\
 &\times
 W_{a/\pi}^{\mathrm{reg}}
 \left(
 \frac{x_D}{y_\pi},
 \kT-\frac{x_D}{y_\pi}\bm p_{\pi T},
 \DT
 \right)
 \cJ_\pi.
 \label{eq:pion-active-cont}
\end{align}
The internal pion parent retains
\begin{verbatim}
VALIDATION_ONLY
PION_PARTON_PARENT_UNMATCHED
LINK_SHORTENING_REQUIRED
UV_MATCHING_REQUIRED
RAPIDITY_SOFT_MATCHING_REQUIRED
NO_EVOLUTION_APPLIED
NO_PROCESS_MAP_APPLIED
\end{verbatim}
Continuum calibration of the nuclear splitting function does not change these
statuses.

\subsection{Transition operators}

For an operator family \(A\), define
\begin{equation}
 \cO^{A}_{QP}(E):\HNN\longrightarrow\Hil_{E\alpha},
 \qquad
 \cO^{A}_{PQ}(E)=\left[\cO^{A}_{QP}(E)\right]^{\dagger}
 \label{eq:transition-operator-cont}
\end{equation}
where supported.  Vector, axial, pseudoscalar, EMT, and partonic transition
operators retain source/target sector, channel, recoil, transfer, helicity,
pion charge, Wilson order, and matching status.

\subsection{Common-parent reductions}

The continuum N2 parent must still satisfy the commuting reductions
\begin{equation}
 W_{a/D}
 \longrightarrow
 \left\{
 \begin{array}{c}
 \TMD\\
 \GPD\\
 \PDF\\
 J^\mu\\
 T^{\mu\nu}
 \end{array}
 \right.
 \label{eq:n2-reductions}
\end{equation}
with one state normalization and one continuum measure.  No named tensor or
pion function may be renormalized independently to repair a failed route.

\section{Separator flow and pion matching}
\label{sec:separator}

\subsection{Resolution separator}

Introduce a separator \(\Lambda_\pi\) and complementary maps
\begin{equation}
 P_{\mathrm{int}}(\Lambda_\pi)
 +P_{\mathrm{exch}}(\Lambda_\pi)
 =I
 \label{eq:separator-partition}
\end{equation}
on the relevant pion-like configuration space, allowing an overlap collar
when the maps are smooth rather than orthogonal.  The exact definition,
regulator, and matching order are part of the separator identity.

\subsection{Matched pion result}

The matched contribution is
\begin{equation}
 W_{\pi}^{\mathrm{matched}}(\Lambda_\pi)
 =
 W_{\pi}^{\mathrm{internal}}(\Lambda_\pi)
 +W_{\pi}^{\mathrm{exchange}}(\Lambda_\pi)
 -W_{\pi}^{\mathrm{overlap}}(\Lambda_\pi).
 \label{eq:matched-pion}
\end{equation}
At a calculation complete through nuclear order \(\nu\), require
\begin{equation}
 \frac{\dd}{\dd\Lambda_\pi}
 W_{\pi}^{\mathrm{matched}}(\Lambda_\pi)
 =\cO(Q^{\nu+1}).
 \label{eq:separator-invariance}
\end{equation}
The same flow applies to the Hamiltonian, current, EMT, partonic operators, and
norm kernel.

\subsection{Flow ledger}

The \SeparatorTrajectory\ records, at each \(\Lambda_\pi\),
\begin{equation}
 \left\{
 W_{\mathrm{int}},
 W_{\mathrm{exch}},
 W_{\mathrm{ov}},
 H_{\ind},
 J_{\ind},
 \cO_{\ind},
 N_{\eff},
 W_{\mathrm{matched}}
 \right\}.
 \label{eq:separator-ledger}
\end{equation}
Large cancellations among components are acceptable only if the matched sum
is stable and the individual numerical errors remain controlled.

\subsection{No-double-counting cells}

The provenance complex contains a two-cell
\begin{equation}
 \texttt{INTERNAL}+\texttt{EXCHANGE}
 \xRightarrow{\ -\texttt{OVERLAP}\ }
 \texttt{MATCHED},
 \label{eq:overlap-two-cell}
\end{equation}
and an explicit/induced equivalence cell.  Missing or duplicate subtraction,
or selection of the explicit continuum together with its fully induced
replacement, fails before numerical evaluation.

\section{Feshbach equivalence with transformed operators}
\label{sec:feshbach}

\subsection{Effective Hamiltonian}

The projected Hamiltonian is
\begin{equation}
 H_{\eff}(E)
 =PHP+PHQ(E-QHQ)^{-1}QHP.
 \label{eq:heff}
\end{equation}
The wave operator is \(\omega(E)=Q(E-QHQ)^{-1}QHP\).

\subsection{Effective operators}

For any operator \(O\), the source-lifted form is
\begin{equation}
 O_{\eff}(E',E)
 =P[1+\omega^{\dagger}(E')]O[1+\omega(E)]P.
 \label{eq:oeff}
\end{equation}
For normalized Hermitian effective-space matrix elements, use
\begin{equation}
 \overline O_{\eff}
 =N(E')^{-1/2}O_{\eff}(E',E)N(E)^{-1/2},
 \label{eq:metric-oeff}
\end{equation}
with the norm kernel from \cref{eq:norm-kernel}.  Applying \(PHP\) while using
\(POP\) is not an equivalent reduction.

\subsection{Operator families required in the comparison}

The explicit/induced benchmark includes:
\begin{itemize}
\item vector charge and current;
\item axial and pseudoscalar operators where supported;
\item EMT operators;
\item pion-active partonic operators;
\item transition operators;
\item the norm kernel and pole-residue map.
\end{itemize}
The relation is
\begin{verbatim}
EXPLICIT_CONTINUUM_NNPI
  EQUIVALENT_TO
INDUCED_NN_OPERATOR_SET
  + TRANSFORMED_CURRENTS
  + TRANSFORMED_PARTONIC_OPERATORS
  + VISIBLE_REMAINDER
\end{verbatim}
never an additive model combination.

\subsection{Unitary representation covariance}

Different off-shell prescriptions or unitary transformations may move strength
among one-body, exchange, contact, and induced terms.  Physical observables
remain invariant only when the Hamiltonian and all operators are transformed
consistently.  This is the same representation-covariance principle seen in
chiral charge operators \cite{Pastore2011}.

\section{Continuum-aware coherent small-\texorpdfstring{\(x\)}{x} pilot}
\label{sec:coherent}

\subsection{Amplitude construction}

The coherent contribution is
\begin{equation}
 \delta W_{\Lambda'\Lambda}^{\mathrm{coh}}
 =
 \sum_{X,\alpha}
 \int\dd\mu_\alpha(E)
 \mathcal A_{\Lambda'\to X,\alpha}^{*}(E)
 \mathcal G_{X,\alpha}(E)
 \mathcal A_{\Lambda\to X,\alpha}(E),
 \label{eq:coherent-continuum}
\end{equation}
with explicit helicities, continuum channels, longitudinal ordering,
propagation phase, species, scattering order, nuclear member, and source.

\subsection{Mandatory zero and reversal limits}

Require
\begin{equation}
 \mathcal A_1=0\quad\text{or}\quad\mathcal A_2=0
 \quad\Longrightarrow\quad
 \delta W^{\mathrm{coh}}=0,
 \label{eq:coherent-zero}
\end{equation}
and the declared phase transformation under longitudinal-order interchange.
Scalar, vector, and tensor responses are projected only after amplitudes are
combined.  A copied unpolarized shadowing ratio must fail the tensor
benchmark.  Physical leading-twist shadowing requires physical diffractive
input and a factorization analysis beyond this pilot \cite{FrankfurtGuzeyStrikman2012}.

\subsection{Parton--nuclear overlap}

The matched result is
\begin{equation}
 W^{\mathrm{matched}}
 =W^{\mathrm{partonic\ Wilson}}
 +W^{\mathrm{nuclear\ coherent}}
 -W^{\mathrm{shared\ soft/Glauber}},
 \label{eq:parton-nuclear-overlap}
\end{equation}
with a count-once two-cell.  The nuclear coherent phase cannot be inserted by
modifying the partonic staple.

\section{Completely positive reductions after coherence}
\label{sec:cp}

\subsection{Amplitude embedding}

Let
\begin{equation}
 V:\Hil_{\mathrm{resolved}}
 \longrightarrow
 \Hil_D\otimes\Hil_{\mathrm{unresolved}}
 \label{eq:amplitude-embedding}
\end{equation}
be the full amplitude after continuum and coherent components are combined.
The reduced map is
\begin{equation}
 \cE(\rho)
 =\Tr_{\mathrm{unresolved}}
 \left[V\rho V^{\dagger}\right].
 \label{eq:cp-map}
\end{equation}
It is completely positive by construction.

\subsection{Premature-trace failure}

Tracing the \(NN\pi\), charge, or coherent labels before the amplitudes are
combined deletes interference terms.  N2 must exhibit at least one current,
tensor, or coherent observable for which
\begin{equation}
 \Tr\left[\sum_\alpha A_\alpha\rho A_\alpha^{\dagger}\right]
 \neq
 \Tr\left[
 \left(\sum_\alpha A_\alpha\right)
 \rho
 \left(\sum_\beta A_\beta\right)^{\dagger}
 \right].
 \label{eq:premature-trace}
\end{equation}
Complete positivity is therefore not a license to discard resolved coherence.

\section{Tagged, pole, and differential diagnostics}
\label{sec:tagged}

\subsection{Tagged continuum channel}

A tagged observable retains spectator momentum and helicity, continuum
energy, charge channel, active constituent, transition status, nuclear
member, and pole variable.  Integrating over the tagged continuum and
spectator variables must recover the inclusive N2 parent.

\subsection{Pole extraction}

Near the active-nucleon pole, the tagged amplitude is decomposed into a pole
term plus a regular remainder.  Multiplication by the appropriate pole factor
must recover the residue implied by the same spectral amplitude and norm
kernel.  Regular final-state terms may change the background but not the pole
residue in the analytic oracle.  Tagged spectator methods provide an
independent test of the light-front spectral function and tensor structure
\cite{CosynWeiss2020}.

\subsection{Pion-resolved closure}

The tagged-to-inclusive test is performed separately for:
\begin{itemize}
\item active proton;
\item active neutron;
\item active pion;
\item resolved \(NN\leftrightarrow NN\pi\) transition;
\item each charge and tensor channel.
\end{itemize}
Detector acceptance remains a separate \(\Proc\)-layer map.

\section{Tensor-network and continuum realization}
\label{sec:tn}

\subsection{State-network topology}

The nuclear state network contains branches
\begin{equation}
 \mathcal N_{\mathrm{state}}
 =
 \mathcal N_{NN}
 \oplus
 \mathcal N_{NN\pi}^{\mathrm{disc}}
 \oplus
 \mathcal N_{NN\pi}^{\mathrm{cont/FV}}.
 \label{eq:n2-network}
\end{equation}
Edges retain sector, charge channel, energy or level index, partial wave,
Jacobi orbitals, pion isospin, nucleon and deuteron helicities, S/D ancestry,
separator, regulator, current/operator channel, microscopic member, and
partonic Wilson identity.

\subsection{Separate truncation axes}

The following are independent axes:
\begin{enumerate}
\item Hamiltonian basis and Fock truncation;
\item TTN bond dimension;
\item continuum level count or volume;
\item continuum quadrature and smearing;
\item separator choice;
\item current/operator order;
\item coherent scattering order;
\item numerical integration.
\end{enumerate}
Convergence in one axis cannot be used to hide another.

\subsection{Observable-sensitive convergence}

Report separately:
\begin{equation}
 \left\{
 E_D,
 Z_D,
 Z_{NN\pi},
 V_{NN\leftrightarrow NN\pi},
 J^\mu,
 \delta_{\mathrm{cont}},
 \delta_{\mathrm{ang}},
 W_{\pi},
 W_{\mathrm{tensor}}^{\mathrm{trans}},
 W_{\mathrm{tensor}}^{\mathrm{coh}},
 W_{\pi}^{\mathrm{matched}}
 \right\}.
 \label{eq:n2-convergence-observables}
\end{equation}
At least one reduced-bond or coarse-continuum approximation must exhibit a
small energy/norm defect but a visible current, transition, or tensor defect.

\section{Assumption compiler and provenance}
\label{sec:assumptions}

\subsection{Primary plans}

The compiler supports at least:
\begin{verbatim}
N2-PLAN-A
  continuum-calibrated NNPI transition
  declared-order complete current basis
  AV18-derived NN branch
  H7 PLAN-A microscopic nucleons

N2-PLAN-B
  finite-volume/discretized-continuum sequence
  same current-generation rules
  Norfolk-derived NN branch
  H7 PLAN-A microscopic nucleons

N2-PLAN-C
  continuum-calibrated transition
  AV18-derived NN branch
  H7 PLAN-B microscopic nucleons

N1-REFERENCE
  immutable C16 analytic-transition parent
\end{verbatim}
These are mutually exclusive complete branches.  The continuum route and its
finite approximation are compared through \FiniteVolumeMap; they are not two
additive mechanisms.

\subsection{Provenance two-complex}

The \ProvComplex\ contains cells for:
\begin{itemize}
\item explicit continuum versus induced operator descriptions;
\item internal/exchange/overlap pion matching;
\item Hamiltonian-to-current gauging;
\item equivalent continuum and finite-volume representations;
\item partonic/nuclear shared soft support;
\item coherent amplitudes and post-coherence CP reduction.
\end{itemize}
A nontrivial unresolved cycle is an audit failure or a genuinely distinct
alternative; it is not interpreted as a physical amplitude.

\section{Uncertainty and covariance}
\label{sec:uncertainty}

The N2 covariance is decomposed into
\begin{equation}
 \Cov_{N2}
 =\Cov_{N}
 +\Cov_{NN}
 +\Cov_{\mathrm{cont}}
 +\Cov_{\mathrm{current}}
 +\Cov_{\mathrm{sep}}
 +\Cov_{\mathrm{coh}}
 +\Cov_{\mathrm{TN}}
 +\Cov_{\num}
 +\Cov_{\mathrm{match}}.
 \label{eq:n2-covariance}
\end{equation}
Here \(\Cov_N\) is the correlated microscopic proton/neutron uncertainty,
\(\Cov_{\mathrm{cont}}\) includes spectral and finite-volume errors,
\(\Cov_{\mathrm{current}}\) is the current-order and low-energy-constant
component, and \(\Cov_{\mathrm{sep}}\) measures residual separator dependence.
These components are not collapsed into one envelope.

A common N2 member propagates the same state, transition, current, separator,
and coherent identities to every observable.  Alternative AV18, Norfolk, and
microscopic assumption plans remain alternative models unless a shared
posterior relates them.

\section{Computational realization and typed interfaces}
\label{sec:software}

\subsection{Core objects}

\begin{longtable}{L{0.28\textwidth}L{0.64\textwidth}}
\toprule
Object & Responsibility\\
\midrule
\endhead
\texttt{N2Resolution} & Nuclear basis, interaction/current order, continuum and separator identity, regulator, and numerical tolerances.\\
\ContinuumChannel & Charge-, spin-, isospin-, partial-wave-, threshold-, and normalization-resolved continuum channel.\\
\SpectralDensity & Continuum measure, support, threshold law, and channel covariance.\\
\FiniteVolumeMap & Discrete levels and weights, boundary conditions, continuum comparison map, and convergence residuals.\\
\TransitionKernel & Hermitian \(NN\leftrightarrow NN\pi\) continuum transition family and generated adjoint.\\
\PoleResidue & Pole condition, principal-value self-energy, cut, derivative, residue, and norm-kernel ledger.\\
\GaugedTerm & One Hamiltonian term with all external-field attachments and proof-backed neutral status.\\
\ExchangeCurrent & Typed one-, two-, transition-, contact-, regulator-, and induced-current block.\\
\shortstack[l]{\texttt{CurrentBasis}\\\texttt{CompletenessCertificate}} & Hamiltonian-term/current incidence map and completeness gate.\\
\BlockContinuity & Sector-, charge-, helicity-, and tensor-resolved signed continuity residuals.\\
\SeparatorTrajectory & Internal, exchange, overlap, induced, and matched flows versus \(\Lambda_\pi\).\\
\shortstack[l]{\texttt{ExplicitInduced}\\\texttt{PionComparison}} & Feshbach Hamiltonian/operator/norm comparison with visible remainder.\\
\ContinuumCoherent & Continuum-channel helicity amplitudes, propagation phases, ordering, and overlap subtraction.\\
\StateBundle & N2 state, continuum map, currents, partonic operators, TTN, ledgers, covariance, and readiness.\\
\PredictionPlan & Frozen assumption branch and executable dependency graph.\\
\ProvComplex & States, operators, transformations, equivalence/subtraction cells, and ancestry.\\
\bottomrule
\end{longtable}

\subsection{Typed data flow}

The executable chain is
\begin{equation}
 \begin{tikzcd}[column sep=small]
 \texttt{N1State}
 \arrow[r]
 &\texttt{ContinuumChannel}
 \arrow[r]
 &\texttt{TransitionKernel}
 \arrow[r]
 &\texttt{PoleResidue}
 \arrow[d]
 \\
 \texttt{Observable}
 &\texttt{ReductionMap}\arrow[l]
 &\texttt{N2Parent}\arrow[l]
 &\texttt{GaugedHamiltonian+Current}\arrow[l]
 \end{tikzcd}
 \label{eq:n2-data-flow}
\end{equation}
Every arrow validates coordinate, channel, member, regulator, separator,
operator, scheme, and provenance identity.

\subsection{Cache identity}

A cache key contains at least
\begin{equation}
 \mathrm{hash}
 \left(
 H_{N2},
 \mu_{\mathrm{cont}},
 V_{PQ},
 J^\mu,
 \Lambda_\pi,
 \cO,
 \text{member},
 \text{kinematics},
 \text{code version}
 \right).
 \label{eq:n2-cache}
\end{equation}
Changing any current attachment, regulator, continuum map, separator, or norm
kernel invalidates the cache.

\section{Formal requirements and property-based tests}
\label{sec:requirements}

\subsection{Stable requirement identifiers}

The normative requirement families are:
\begin{longtable}{L{0.20\textwidth}L{0.72\textwidth}}
\toprule
ID & Requirement\\
\midrule
\endhead
\texttt{V14.CONT.1--6} & Direct-integral channels, normalization, support, thresholds, and projectors.\\
\texttt{V14.FV.1--6} & Finite spectral maps, principal-value convergence, and terminology gates.\\
\texttt{V14.POLE.1--7} & Self-energy, cut, pole, residue, norm kernel, and no-width theorem.\\
\texttt{V14.CAL.1--4} & Restricted calibration, frozen holdouts, Jacobian, and null directions.\\
\texttt{V14.GAUGE.1--7} & Gauged Hamiltonian, nonlocal kernels, regulator currents, and term ownership.\\
\texttt{V14.CURR.1--8} & Declared-order current basis and completeness certificate.\\
\texttt{V14.CONTI.1--8} & Blockwise continuity, charge/tensor resolution, current moments, and angular condition.\\
\texttt{V14.PION.1--6} & Continuum pion-active and transition operators with unmatched status.\\
\texttt{V14.SEP.1--7} & Separator trajectory, stable matched total, and no-double-counting cells.\\
\texttt{V14.FESH.1--6} & Hamiltonian/operator/norm transformation and representation covariance.\\
\texttt{V14.COH.1--6} & Continuum helicity amplitudes, zero/reversal limits, and shared-soft subtraction.\\
\texttt{V14.CP.1--4} & Post-coherence CP map and premature-trace failure.\\
\texttt{V14.TN.1--6} & Continuum TTN, full-bond equality, and observable-sensitive convergence.\\
\texttt{V14.PROV.1--5} & Assumption plans, two-complex normalization, and ancestry.\\
\texttt{V14.GATE.1--8} & Downstream matching, sector, process, inference, and production gates.\\
\bottomrule
\end{longtable}

\subsection{Core property tests}

The implementation must test at least:
\begin{enumerate}
\item continuum orthogonality and resolution of identity;
\item exact zero discontinuity below threshold;
\item principal-value and cut signs;
\item finite-volume-to-continuum convergence;
\item transition-kernel Hermiticity;
\item pole position, residue, and norm-kernel agreement;
\item one current attachment certificate per Hamiltonian term;
\item blockwise and charge-resolved continuity;
\item direct-current versus GTMD/GPD moments;
\item separator stability and signed missing/duplicate-subtraction defects;
\item explicit/induced equivalence with transformed operators and remainder;
\item coherent zero, order-reversal, tensor, and overlap tests;
\item CP/Kraus equality after coherence and premature-trace failure;
\item exact/Krylov/full-bond agreement;
\item distinct TTN, continuum, separator, and numerical convergence axes.
\end{enumerate}

\subsection{Required negative-test scale}

A production implementation should contain at least 340 stable negative-test
identities covering continuum support, transition dynamics, current
completeness, blockwise continuity, pion separation, coherent/CP ordering,
tensor-network compression, stale caches, and all downstream readiness gates.
Every failure carries a stable ID, expected exception class, and deterministic
diagnostic.

\section{Analytic and finite-dimensional benchmark suite}
\label{sec:benchmarks}

\subsection{N2-A: analytic separable continuum}

Choose
\begin{equation}
 V_\alpha(E)=g_\alpha
 (E-E_{\mathrm{th},\alpha})^{\nu_\alpha/2}
 e^{-(E-E_{\mathrm{th},\alpha})/(2\Lambda_\alpha)}
 \theta(E-E_{\mathrm{th},\alpha}),
 \label{eq:separable-oracle}
\end{equation}
and evaluate the principal value, cut, pole, derivative, and residue by an
independent quadrature or analytic expression.  Below-threshold absorption
must be exactly zero.

\subsection{N2-B: finite-volume convergence}

Construct a sequence of levels and weights converging to
\cref{eq:separable-oracle}.  Test the principal value, cut, pole, residue,
transition matrix element, current matrix element, and pion-active moment.

\subsection{N2-C: charge-complete continuum channels}

Verify charge \(+1\), isospin zero, parity, normalization, and channel
orthogonality for \(pn\pi^0\), \(pp\pi^-\), and \(nn\pi^+\).

\subsection{N2-D: transition-kernel Hermiticity}

Check emission/absorption adjoints on basis states and deterministic complex
superpositions, including regulator, recoil, phase, and charge-channel
identity.

\subsection{N2-E: Hamiltonian gauging}

Generate currents from a finite Hamiltonian containing kinetic,
sector-changing, derivative, contact, nonlocal-regulator, and induced terms.
Every retained term must appear in the completeness certificate.

\subsection{N2-F: blockwise continuity}

Verify \cref{eq:continuity-pp,eq:continuity-qp} and the conjugate blocks in
every charge channel.  Removing each nonzero current contribution must produce
a signed defect.

\subsection{N2-G: current-component and angular closure}

Construct a covariant spin-1 current, generate all supported light-front
helicity amplitudes, and compare charge, magnetic, quadrupole, angular, and
GTMD/GPD-moment routes.

\subsection{N2-H: separator flow}

Vary \(\Lambda_\pi\) so that internal, exchange, overlap, and induced terms
move substantially while the matched total remains stable within the declared
order.

\subsection{N2-I: Feshbach Hamiltonian/operator equivalence}

Compare full and effective calculations for the Hamiltonian, vector current,
axial/pseudoscalar operator, EMT, pion-active partonic operator, transition
operator, and norm kernel.  Omitting an induced operator must fail.

\subsection{N2-J: pion-active common-parent closure}

Verify direct and sequential TMD/GPD/PDF/current reductions from the
continuum-calibrated nuclear splitting amplitude without a fitted pion
normalization.

\subsection{N2-K: transition tensor signal}

Construct one tensor observable that requires a specific transition partial
wave and vanishes when either participating amplitude is removed.

\subsection{N2-L: coherent continuum pilot}

Test zero elementary amplitudes, longitudinal-order reversal, scalar/vector/
tensor projections, failure of an unpolarized ratio, and shared-soft
subtraction.

\subsection{N2-M: CP reduction}

Compare the explicit post-coherence partial trace with the Kraus form and show
that an early trace loses an interference observable.

\subsection{N2-N: exact, Krylov, TTN, and continuum convergence}

Full bond and fine continuum must reproduce the exact oracle.  At least one
reduced-bond or coarse-continuum state must retain a small energy error while
losing a real transition/current/tensor signal.

\subsection{N2-O: tagged-inclusive closure}

Integrate continuum- and pion-resolved tagged observables and recover the
inclusive N2 parent and the declared pole residue.

\subsection{N2-P: frozen holdouts and provenance normalization}

Keep at least one transition, residue, pion, tensor, current, angular, and
coherent observable unfitted.  Verify that every explicit/induced, overlap,
current, coherent, and CP path is counted exactly once.

\section{Staged N2 implementation program}
\label{sec:stages}

\subsection{N2.0: continuum and finite-volume types}

Implement \ContinuumChannel, \SpectralDensity, \FiniteVolumeMap, support,
normalization, thresholds, and analytic benchmarks.

\subsection{N2.1: transition, pole, and residue}

Implement \TransitionKernel, self-energy, principal value, cut, pole,
derivative, residue, norm kernel, and calibration/holdout manifests.

\subsection{N2.2: gauged Hamiltonian and currents}

Implement the gauged-Hamiltonian-term, exchange-current, and current-basis
completeness interfaces.  Add regulator gauging and blockwise continuity.

\subsection{N2.3: pion operators and separator flow}

Implement continuum-aware pion-active and transition operators,
\SeparatorTrajectory, Feshbach comparisons, and deterministic two-cells.

\subsection{N2.4: coherent, CP, and tagged pilots}

Upgrade the coherent amplitude, shared-soft map, CP reduction, and tagged
continuum diagnostics.

\subsection{N2.5: tensor network, convergence, and release gates}

Integrate continuum indices into the TTN, complete all convergence manifests,
negative tests, assumption plans, cache validation, and fail-closed downstream
gates.

\section{Acceptance criteria and scientific status}
\label{sec:acceptance}

N2 is formally accepted only when:
\begin{enumerate}
\item continuum and finite-volume routes have typed normalization and support;
\item the finite sequence converges to the analytic continuum oracle;
\item pole, residue, principal-value, cut, and norm ledgers close;
\item below-threshold absorption is exactly zero and numerical \(\eps\) is not physical;
\item the transition kernel is Hermitian with generated adjoints;
\item every retained Hamiltonian term has certified current attachments or a proof-backed neutral status;
\item the declared-order current basis has no unexplained gap;
\item continuity closes separately in every sector, charge, helicity, and tensor block;
\item charge, magnetic, quadrupole, current-component, angular, and moment diagnostics pass;
\item the continuum pion-active route closes and remains explicitly unmatched;
\item the matched pion total is stable under separator variation;
\item Feshbach elimination transforms Hamiltonian, currents, partonic operators, and norm kernel with a visible remainder;
\item the coherent pilot is helicity resolved and explicitly nonphysical;
\item partonic Wilson and nuclear coherent mechanisms remain separate with one overlap subtraction;
\item CP reduction occurs only after amplitudes are combined;
\item exact, Krylov, TTN, finite-volume, separator, and numerical convergence are reported separately;
\item nontrivial transition, residue, pion, tensor, current, and coherent holdouts remain unfitted;
\item all negative tests and prior regression gates pass;
\item production artifacts and registries remain unchanged;
\item all manifests rebuild deterministically;
\item all statuses remain qualified and validation only;
\item \(\Delta\Delta\), compact six-quark, hidden-color, physical matching, evolution, process, and inference gates remain closed.
\end{enumerate}

Allowed statuses include:
\begin{verbatim}
N2_CONTINUUM_NNPI_TRANSITION_VALIDATED
N2_FINITE_VOLUME_TO_CONTINUUM_MAP_VALIDATED
N2_POLE_AND_RESIDUE_BENCHMARKED
N2_DECLARED_ORDER_EXCHANGE_CURRENT_BASIS_COMPLETE
N2_FINITE_BASIS_CONTINUITY_CLOSED
N2_SEPARATOR_STABILITY_VALIDATED
N2_PION_ACTIVE_ROUTE_VALIDATED_UNMATCHED
N2_COHERENT_CONTINUUM_PILOT_VALIDATED
N2_CP_REDUCTION_VALIDATED
N2_TTN_CONTINUUM_BRANCH_VALIDATED
N2_VALIDATION_ONLY
\end{verbatim}

\section{Boundary to N3 and later QCD matching}
\label{sec:boundary}

If the continuum-transition, current-completeness, separator, and coherent
checks close, the next nuclear layer may add
\begin{equation}
 \Hil_{\mathrm{N3}}
 =\Hil_{\mathrm{N2}}
 \oplus\Hil_{\Delta\Delta}
 \oplus\Hil_{6q}
 \oplus\cdots,
 \label{eq:n3-space}
\end{equation}
with explicit transition currents, normalized interference, color and spin
structure, and matched coherent amplitudes.  A scalar \(\Delta\Delta\) or
six-quark probability remains insufficient.

Independently, the complete low-resolution nuclear parent still requires the
LF-to-QCD matching, soft/rapidity completion, rank-aware evolution, process
factorization, and inference layers.  Continuum calibration of the nuclear
pion sector does not close those gates.

\section{Conclusion}

Volume XIV completes the formal N1-to-N2 transition.  The pion sector is no
longer merely a finite analytic branch: it is a direct-integral continuum with
a controlled finite spectral approximation, a dispersive self-energy, a
physical support rule, a pole and residue, and a normalized operator metric.
The nuclear current is no longer a curated list detached from the force:
every retained Hamiltonian term is gauged, every attachment is certified, and
continuity is tested block by block and channel by channel.  The pion
separator moves strength among internal, exchange, overlap, and induced
representations while the matched total remains stable.  Currents, EMT
operators, pion-active and transition partonic operators, and the norm kernel
transform with the Hamiltonian.  Coherent amplitudes are combined before
partial trace and remain distinct from the partonic Wilson staple.  This is
the required foundation for adding explicit \(\Delta\Delta\), compact
six-quark, and hidden-color sectors without rebuilding the nuclear model as a
patchwork of separately normalized corrections.

\appendix

\section{Derivation of the self-energy discontinuity}
\label{app:disc}

Using
\begin{equation}
 \frac{1}{x+i0}=\PV\frac1x-i\pi\delta(x),
 \label{eq:distribution-identity}
\end{equation}
insert the spectral projector \cref{eq:spectral-projector} into
\cref{eq:projected-resolvent}.  For a matrix element between \(P\)-space
states \(|a\rangle\), \(|b\rangle\),
\begin{equation}
 \Sigma_{ab}(E_0+i0)
 =\sum_\alpha\int\dd E\,
 \rho_\alpha(E)
 \frac{V_{a\alpha}(E)V_{\alpha b}(E)}{E_0-E+i0}.
 \label{eq:self-energy-app}
\end{equation}
Applying \cref{eq:distribution-identity} gives \cref{eq:self-energy-disc}.
The delta term contributes only when \(E_0\) lies inside the support of
\(\rho_\alpha\).  A finite numerical width approximates this distribution but
does not define the support.

\section{Pole-residue and norm-kernel relation}
\label{app:residue}

For a scalar projected resolvent
\begin{equation}
 G_P(E)=\frac{1}{E-E_0-\Sigma(E)},
\end{equation}
expand around a simple pole \(E_D\):
\begin{equation}
 E-E_0-\Sigma(E)
 =
 (E-E_D)\left[1-\Sigma'(E_D)\right]+ \cO((E-E_D)^2).
\end{equation}
Thus
\begin{equation}
 G_P(E)
 =\frac{Z_D}{E-E_D}+\text{regular},
 \qquad
 Z_D=\frac{1}{1-\Sigma'(E_D)}.
\end{equation}
The same derivative appears in the normalization of the projected component.
In a matrix problem, \(Z_D\) becomes the residue on the pole eigenspace and
must be evaluated with the proper metric.

\section{Block continuity derivation}
\label{app:block-continuity}

Multiplying the block matrices in \cref{eq:block-operators} yields
\begin{align}
 [H,\rho]_{PP}
 &=H_P\rho_{PP}+V_{PQ}\rho_{QP}
 -\rho_{PP}H_P-\rho_{PQ}V_{QP},\\
 [H,\rho]_{PQ}
 &=H_P\rho_{PQ}+V_{PQ}\rho_{QQ}
 -\rho_{PP}V_{PQ}-\rho_{PQ}H_Q,\\
 [H,\rho]_{QP}
 &=V_{QP}\rho_{PP}+H_Q\rho_{QP}
 -\rho_{QP}H_P-\rho_{QQ}V_{QP},\\
 [H,\rho]_{QQ}
 &=V_{QP}\rho_{PQ}+H_Q\rho_{QQ}
 -\rho_{QP}V_{PQ}-\rho_{QQ}H_Q.
\end{align}
Each block must be matched by the divergence of the corresponding current
block.  A diagonal one-body current cannot close the transition identities.

\section{Machine-readable schemas}
\label{app:schemas}

\subsection{Continuum-channel record}

\begin{verbatim}
{
  "channel_id": "...",
  "charge_channel": "pn_pi0",
  "total_charge": 1,
  "isospin": 0,
  "J_parity": "1+",
  "Jz": "...",
  "pair_spin": "...",
  "pair_isospin": "...",
  "jacobi_orbitals": ["...", "..."],
  "threshold": "...",
  "normalization": "ENERGY_DELTA",
  "spectral_density_id": "...",
  "regulator_id": "...",
  "member_id": "..."
}
\end{verbatim}

\subsection{Pole and residue record}

\begin{verbatim}
{
  "pole_id": "...",
  "hamiltonian_id": "...",
  "transition_kernel_id": "...",
  "spectral_map_id": "...",
  "pole_energy": "...",
  "principal_value": "...",
  "cut_weight": "...",
  "self_energy_derivative": "...",
  "residue": "...",
  "norm_kernel_id": "...",
  "below_threshold_absorption": 0.0,
  "epsilon_is_physical": false
}
\end{verbatim}

\subsection{Current completeness certificate}

\begin{verbatim}
{
  "certificate_id": "...",
  "hamiltonian_id": "...",
  "declared_order": "...",
  "terms": [
    {
      "hamiltonian_term_id": "...",
      "charge_status": "CHARGED",
      "current_attachments": ["..."],
      "regulator_attachment": "...",
      "contact_partners": ["..."],
      "continuity_blocks": ["NN_TO_NN", "NN_TO_NNPI"],
      "unresolved_gaps": []
    }
  ],
  "complete_at_declared_order": true
}
\end{verbatim}

\subsection{Separator trajectory record}

\begin{verbatim}
{
  "trajectory_id": "...",
  "separator_values": ["..."],
  "internal": ["..."],
  "exchange": ["..."],
  "overlap": ["..."],
  "induced_hamiltonian": ["..."],
  "induced_current": ["..."],
  "induced_partonic_operator": ["..."],
  "matched_total": ["..."],
  "stability_residual": "...",
  "truncation_tolerance": "..."
}
\end{verbatim}

\section{Compact N2 acceptance checklist}

\begin{enumerate}
\item Direct-integral continuum channels normalized and support typed?
\item Complete charge, isospin, parity, helicity, and partial-wave identity?
\item Finite-volume/discrete sequence converges to the continuum oracle?
\item Principal value, cut, pole, derivative, residue, and norm kernel close?
\item Exact zero absorption below threshold and no physical numerical \(\eps\)?
\item Hermitian transition kernel with generated adjoint?
\item Restricted calibration and frozen holdouts?
\item Every Hamiltonian term gauged or certified neutral?
\item Complete declared-order current basis with no unexplained gap?
\item Blockwise continuity in all sectors, charges, helicities, and tensor blocks?
\item Current-component, angular, and GTMD/GPD-moment closure?
\item Continuum pion-active and transition routes close while unmatched?
\item Stable separator trajectory for Hamiltonian and all operators?
\item Explicit/induced Feshbach equivalence with metric and visible remainder?
\item Helicity-resolved coherent pilot with exact zero/reversal limits?
\item Partonic/nuclear shared-soft subtraction counted once?
\item CP map derived only after amplitudes are combined?
\item Tagged-inclusive and pole-residue closure?
\item Exact/Krylov/full-bond equality and observable-sensitive convergence?
\item Deterministic assumption plans, provenance, caches, and manifests?
\item All unavailable sectors and downstream gates fail closed?
\item Production and authoritative artifacts unchanged?
\end{enumerate}

\small
\begin{thebibliography}{99}

\bibitem{Dirac1949}
P.~A.~M.~Dirac,
``Forms of relativistic dynamics,''
Rev.\ Mod.\ Phys.\ \textbf{21}, 392 (1949).

\bibitem{KeisterPolyzou1991}
B.~D.~Keister and W.~N.~Polyzou,
``Relativistic Hamiltonian dynamics in nuclear and particle physics,''
Adv.\ Nucl.\ Phys.\ \textbf{20}, 225 (1991).

\bibitem{Carbonell1998}
J.~Carbonell, B.~Desplanques, V.~A.~Karmanov, and J.-F.~Mathiot,
``Explicitly covariant light-front dynamics and relativistic few-body systems,''
Phys.\ Rept.\ \textbf{300}, 215 (1998), arXiv:nucl-th/9804029.

\bibitem{Feshbach1962}
H.~Feshbach,
``A unified theory of nuclear reactions. II,''
Ann.\ Phys.\ \textbf{19}, 287 (1962).

\bibitem{KvinikhidzeBlankleider1999I}
A.~N.~Kvinikhidze and B.~Blankleider,
``Gauging of equations method. I. Electromagnetic currents of three
distinguishable particles,''
Phys.\ Rev.\ C \textbf{60}, 044003 (1999), arXiv:nucl-th/9901001.

\bibitem{KvinikhidzeBlankleider1999II}
A.~N.~Kvinikhidze and B.~Blankleider,
``Gauging of equations method. II. Electromagnetic currents of three identical
particles,''
Phys.\ Rev.\ C \textbf{60}, 044004 (1999), arXiv:nucl-th/9901002.

\bibitem{Pastore2009}
S.~Pastore, L.~Girlanda, R.~Schiavilla, M.~Viviani, and R.~B.~Wiringa,
``Electromagnetic currents and magnetic moments in chiral effective field
theory,''
Phys.\ Rev.\ C \textbf{80}, 034004 (2009), arXiv:0906.1800.

\bibitem{Koelling2009}
S.~K\"olling, E.~Epelbaum, H.~Krebs, and U.-G.~Mei\ss ner,
``Two-pion exchange electromagnetic current in chiral effective field theory
using the method of unitary transformation,''
Phys.\ Rev.\ C \textbf{80}, 045502 (2009), arXiv:0907.3437.

\bibitem{Koelling2011}
S.~K\"olling, E.~Epelbaum, H.~Krebs, and U.-G.~Mei\ss ner,
``Two-nucleon electromagnetic current in chiral effective field theory:
one-pion exchange and short-range contributions,''
Phys.\ Rev.\ C \textbf{84}, 054008 (2011), arXiv:1107.0602.

\bibitem{Pastore2011}
S.~Pastore, L.~Girlanda, R.~Schiavilla, and M.~Viviani,
``The two-nucleon electromagnetic charge operator in chiral effective field
theory up to one loop,''
Phys.\ Rev.\ C \textbf{84}, 024001 (2011), arXiv:1106.4539.

\bibitem{Krebs2020}
H.~Krebs,
``Nuclear currents in chiral effective field theory,''
Eur.\ Phys.\ J.\ A \textbf{56}, 234 (2020), arXiv:2008.00974.

\bibitem{Baroni2021}
A.~Baroni, G.~B.~King, and S.~Pastore,
``Electroweak currents from chiral effective field theory,''
Few-Body Syst.\ \textbf{62}, 114 (2021), arXiv:2107.10721.

\bibitem{HeZahed2024}
F.~He and I.~Zahed,
``Deuteron gravitational form factors: exchange currents,''
Phys.\ Rev.\ C \textbf{110}, 014312 (2024), arXiv:2401.09318.

\bibitem{LellouchLuscher2001}
L.~Lellouch and M.~L\"uscher,
``Weak transition matrix elements from finite-volume correlation functions,''
Commun.\ Math.\ Phys.\ \textbf{219}, 31 (2001), arXiv:hep-lat/0003023.

\bibitem{Wiringa1995}
R.~B.~Wiringa, V.~G.~J.~Stoks, and R.~Schiavilla,
``Accurate nucleon-nucleon potential with charge-independence breaking,''
Phys.\ Rev.\ C \textbf{51}, 38 (1995), arXiv:nucl-th/9408016.

\bibitem{Machleidt2001}
R.~Machleidt,
``The high-precision, charge-dependent Bonn nucleon-nucleon potential,''
Phys.\ Rev.\ C \textbf{63}, 024001 (2001), arXiv:nucl-th/0006014.

\bibitem{CaoEtAl2026}
X.~Cao, S.~Cheng, Y.~Duan, Y.~Li, S.~Xu, and X.~Zhao,
``Non-perturbative flavor asymmetry in the nucleon and deuteron: The
light-front Hamiltonian effective field theory approach,''
arXiv:2601.13567 (2026).

\bibitem{FrankfurtGuzeyStrikman2012}
L.~Frankfurt, V.~Guzey, and M.~Strikman,
``Leading twist nuclear shadowing phenomena in hard processes with nuclei,''
Phys.\ Rept.\ \textbf{512}, 255 (2012), arXiv:1106.2091.

\bibitem{CosynWeiss2020}
W.~Cosyn and C.~Weiss,
``Polarized electron--deuteron deep-inelastic scattering with spectator
nucleon tagging,''
Phys.\ Rev.\ C \textbf{102}, 065204 (2020), arXiv:2006.03033.

\bibitem{Miller2014}
G.~A.~Miller,
``Pionic and hidden-color, six-quark contributions to the deuteron \(b_1\)
structure function,''
Phys.\ Rev.\ C \textbf{89}, 045203 (2014), arXiv:1311.4561.

\bibitem{HoodbhoyJaffeManohar1989}
P.~Hoodbhoy, R.~L.~Jaffe, and A.~Manohar,
``Novel effects in deep-inelastic scattering from spin-one hadrons,''
Nucl.\ Phys.\ B \textbf{312}, 571 (1989).

\bibitem{Collins2011}
J.~C.~Collins,
\emph{Foundations of Perturbative QCD},
Cambridge University Press (2011).

\end{thebibliography}

\end{document}
